\documentclass[11pt]{article}
\usepackage[margin=1in]{geometry}
\usepackage{longtable}
\usepackage{booktabs}
\usepackage{xcolor}
\usepackage{hyperref}
\usepackage{listings}
\usepackage{parskip}
\usepackage{enumitem}
\lstset{basicstyle=\ttfamily\small,breaklines=true,frame=single}
\title{BVBRC-111 --- Independent Replication of Harmer \emph{et al.} 2022:\\
       Complete genome of XDR GC1 \emph{Acinetobacter baumannii} MRSN 56\\
       and a novel route to fluoroquinolone resistance}
\author{OpenClaw replication team (Ollie subagent)}
\date{2026-07-05}
\begin{document}
\maketitle

\section{Paper summary}

Harmer, Lebreton, Stam, McGann, Hall.
\emph{J Antimicrob Chemother} 77(7):1851--1855, 2022.
DOI \texttt{10.1093/jac/dkac115}; PMID 35403193; PMCID PMC9244215.

MRSN~56 is an extensively drug-resistant (XDR) global-clone-1 (GC1)
\emph{A.~baumannii} isolate from a US military treatment facility.
The authors combined Oxford Nanopore MinION long reads with existing
Illumina MiSeq short reads through Unicycler to produce a finished
genome (one chromosome $+$ four small plasmids).  Their principal
findings are:

\begin{itemize}[nosep]
  \item ST1 in the Pasteur scheme (typing).
  \item Complete architecture: chromosome plus four plasmids, none carrying
        acquired AMR genes.
  \item An AbaR28 island in \emph{comM} carrying \texttt{aphA1},
        \texttt{aacC1}, \texttt{aadA1} and \texttt{sul1}.
  \item Two copies of Tn\emph{2006} (each carrying \emph{bla}$_{\text{OXA-23}}$).
  \item Two copies of Tn\emph{7} carrying \texttt{dfrA1}, \texttt{sat},
        \texttt{aadA1}.
  \item A novel ``Tn\emph{7}$^{+}$'' configuration: canonical Tn\emph{7}
        at the attTn\emph{7} site downstream of \emph{glmS} plus an
        adjacent AbGRI1-derived $\sim$22.85\,kb segment containing
        \texttt{tet(B)} and \texttt{sul2}.
  \item A novel two-step route to fluoroquinolone resistance =
        a QRDR mutation in \emph{gyrA} (\emph{S83L}) combined with
        IS\emph{Aba1}-mediated inactivation of the \emph{marR}
        repressor, allowing constitutive \emph{marA} expression from an
        IS\emph{Aba1}-internal promoter.
\end{itemize}

\section{Claims table}

\begin{longtable}{@{}lllll@{}}
\toprule
ID & Type & Claim (short) & Testable? & Tested? \\ \midrule
C1 & Typing        & ST1(Pasteur), ST231(Ox), KL1, OCL1                                     & Yes & Yes \\
C2 & Architecture  & Chromosome + 4 plasmids; no AMR on any plasmid                         & Yes & Yes \\
C3 & AMR arch.     & AbaR28 in \emph{comM} carries aphA1/aacC1/aadA1/sul1                   & Yes & Yes \\
C4 & AMR arch.     & Two Tn\emph{2006} (\emph{bla}$_{\text{OXA-23}}$) copies                & Yes & Yes \\
C5 & AMR arch.     & Two Tn\emph{7} copies with dfrA1/aadA1/sat                             & Yes & Yes \\
C6 & Novel arch.   & Tn\emph{7}$^{+}$ downstream of \emph{glmS} with tet(B)/sul2, $\sim$22.85\,kb & Yes & Yes \\
C7 & Mechanism     & gyrA S83L + IS\emph{Aba1}-inactivated \emph{marR} $\to$ constitutive \emph{marA} &
                                                                                             Struct.~yes, Trans.~no &
                                                                                             Struct.~yes, Trans.~no \\
\bottomrule
\end{longtable}

\section{Method}

\subsection{Data acquisition}
All sequences fetched from NCBI on 2026-07-05.
\begin{enumerate}[nosep]
  \item Assembly \texttt{GCA\_021484925.1\_ASM2148492v1}: chromosome
        CP090606.1, plus per-record protein/GFF/GBFF/feature-table.
  \item Plasmids CP080453--CP080456 via \texttt{efetch}.
  \item Reference WT GyrA \texttt{WP\_000116449.1} for QRDR comparison.
  \item Concatenated FASTA \texttt{MRSN56\_complete.fna} = 4{,}174{,}182\,bp.
\end{enumerate}

\subsection{Compute environment}
Heavy compute on \texttt{uicgpu} (8$\times$A100, 255 cores, 2\,TB RAM),
in the pre-existing \texttt{/data/stevens/envs/bvbrc14} environment.

\begin{itemize}[nosep]
  \item \texttt{abricate} 1.4.0 with April-2026 database snapshots
        (resfinder 3206, card 6052, ncbi 8232, megares 6635,
        argannot 2224, plasmidfinder 488, vfdb 4592).
  \item \texttt{mlst} 2.33.1 (schemes \texttt{abaumannii\_2} Pasteur and
        \texttt{abaumannii} Oxford).
  \item \texttt{BLAST+} 2.16.
  \item \texttt{Biopython} 1.85 (feature parsing, translation, alignment).
  \item Argo proxy at \texttt{localhost:44497} for LLM judging:
        \texttt{argo:gpt-5.1} and \texttt{argo:gemini-2.5-pro}.
        No paid API was used.
\end{itemize}

\subsection{Analyses}
\begin{enumerate}[nosep]
  \item Per-contig lengths and total bp.
  \item Both Pasteur and Oxford MLST schemes.
  \item Whole-genome AMR detection against ResFinder / CARD / NCBI /
        MEGARES / ARG-ANNOT at \texttt{--minid 90 --mincov 80}.
  \item Chromosome-only AMR detection (to show plasmids are AMR-free).
  \item Per-plasmid AMR detection (CP080453--CP080456 individually).
  \item PlasmidFinder replicon typing.
  \item VFDB virulence-factor screen.
  \item Feature-level analysis of Tn\emph{7} machinery
        (TnsA/B/C/D/E) and their chromosomal loci.
  \item Copy count and position of IS\emph{Aba1} and IS\emph{26}.
  \item AbaR28-locus walk (chr 340k--410k) CDS by CDS.
  \item IS\emph{Aba1}$\leftrightarrow$\emph{marR} zoom (chr 2{,}310k--2{,}325k).
  \item GyrA QRDR verification: extract CDS, translate, align to
        WP\_000116449.1, enumerate every amino-acid difference.
\end{enumerate}

\subsection{LLM judging}
Both judges were fed the same structured evidence dossier and asked to
score 0--100 with an explicit verdict.
\begin{itemize}[nosep]
  \item \texttt{argo:gpt-5.1} $\to$ 88 / REPLICATED.
  \item \texttt{argo:gemini-2.5-pro} $\to$ 95 / FULLY REPLICATED.
  \item Mean 91.5, rounded verdict REPLICATED (92).
\end{itemize}

\section{Results vs.\ paper}

\subsection{MLST (C1)}
\begin{lstlisting}
mlst --scheme abaumannii_2 GCA_021484925.1_ASM2148492v1_genomic.fna -> abaumannii_2 ST1
  Pas_cpn60(1) Pas_fusA(1) Pas_gltA(1) Pas_pyrG(1)
  Pas_recA(5) Pas_rplB(1) Pas_rpoB(1)
\end{lstlisting}
Paper claims ST1(Pasteur); we recover ST1(Pasteur).  \textbf{MATCH.}
Oxford scheme returned ``$-$'' (novel combination) because the
\texttt{Oxf\_gdhB} allele resolved as 4/182 (mixed), reflecting
pubMLST database drift since 2022 rather than a genome difference.

\subsection{Plasmids carry no AMR (C2)}
\begin{lstlisting}
CP080453 (2178 bp): resfinder=0 card=0 ncbi=0
CP080454 (2725 bp): resfinder=0 card=0 ncbi=0
CP080455 (6772 bp): resfinder=0 card=0 ncbi=0
CP080456 (8731 bp): resfinder=0 card=0 ncbi=0

chromosome-only == whole-genome:
  resfinder 70 == 70; card 93 == 93; ncbi 73 == 73.
\end{lstlisting}
Every acquired AMR gene is chromosomal.  \textbf{MATCH.}

\subsection{AbaR28 in \emph{comM} (C3)}
Chromosomal walk 374{,}829..380{,}957 recovers \texttt{intI1},
\texttt{aac(3)-Ia}, \texttt{aadA1} (pseudo), \texttt{qacE$\Delta$1},
\texttt{sul1} in the canonical class-1 integron order.
\texttt{aphA1} = \texttt{aph(3')-Ia} is present in $\sim$40 copies as
an IS\emph{26}-flanked Tn\emph{6020} tandem amplification at
chr 139k--260k (repeat unit $\sim$2.25\,kb).
All four AbaR28 markers accounted for.  \textbf{MATCH.}

\subsection{Two Tn\emph{2006} / bla$_{\text{OXA-23}}$ (C4)}
Both \emph{bla}$_{\text{OXA-23}}$ copies are flanked by IS\emph{Aba1}
on both sides at chr 1.08\,Mb and chr 4.02\,Mb (Tn\emph{2006}
signature).  \textbf{MATCH.}

\subsection{Two Tn\emph{7} copies + Tn\emph{7}$^{+}$ (C5, C6)}
\begin{itemize}[nosep]
  \item Primary site: \emph{glmS} 4{,}030{,}861..4{,}032{,}700 followed
        immediately by TnsA at 4{,}032{,}858; TnsD; TnsE; \texttt{aadA1};
        \texttt{dfrA1}; \texttt{aph(6)-Id}; \texttt{tet(B)};
        \texttt{sul2}; right-flanking IS\emph{Aba1}.
        Tn\emph{7}$^{+}$ span $\sim$21.7\,kb versus 22{,}852\,bp
        reported by the authors (same order of magnitude, boundary-definition
        difference).
  \item Secondary site: TnsE, TnsD, TnsA at chr 2{,}238k--2{,}246k with
        \texttt{dfrA1}$+$\texttt{aadA1}.
  \item \texttt{sat} recovered as SAT-2 (CARD) / \texttt{sat2\_fam}
        (AMRFinder).
\end{itemize}
\textbf{MATCH.}

\subsection{gyrA QRDR + IS\emph{Aba1}/\emph{marR} (C7)}
\texttt{gyrA} CDS at chr 3{,}203{,}782..3{,}206{,}497 ($-$), 905\,aa;
aligned against WT \texttt{WP\_000116449.1}: 902/904 identity;
\emph{S81L} (i.e.\ classical S83L in older numbering) in the QRDR and
one non-QRDR substitution.
Locus zoom chr 2{,}310k--2{,}320k: IS\emph{Aba1} transposase (pseudo)
immediately upstream of a MarR-family regulator; downstream MFS
transporter.  A second MarR-family regulator at chr 1{,}422k is a
PGAP-flagged frameshifted pseudogene.  20 IS\emph{Aba1} copies
genome-wide provide ample substrate for insertion-mediated inactivation.
Structural prerequisite confirmed; transcriptional activation
untested.  \textbf{MATCH-structural, PARTIAL for transcriptional.}

\subsection{Bonus finding}
\emph{parC} pseudogene at chr 361{,}347..363{,}566 (frameshifted, PGAP).
An additional FQR-associated locus disruption not spelled out in the
paper's abstract.

\section{Per-claim: what worked, what didn't}

\begin{longtable}{@{}p{0.5cm}p{5cm}p{9cm}@{}}
\toprule
ID & Worked & Didn't work / limitations \\ \midrule
C1 & Exact Pasteur ST1 recovered. & Oxford scheme returned ``$-$'';
      diagnosed as pubMLST drift, not a real disagreement. \\
C2 & Zero AMR on all four plasmids across three DBs; identity chromosome-only vs.\ whole-genome. &
     PlasmidFinder identified no replicons on the small plasmids (expected --- they may lack
     canonical rep gene families) so we could not double-check plasmid identity independently. \\
C3 & All four markers found; canonical integron order intact. &
     \texttt{aphA1} is not in the integron proper; it is an IS\emph{26}-mediated Tn\emph{6020}
     tandem amplification (paper's companion manuscript describes this).  Our claim of ``all four''
     therefore lumps in-integron with adjacent IS\emph{26} amplification. \\
C4 & Two \emph{bla}$_{\text{OXA-23}}$ copies with symmetric IS\emph{Aba1} flanks; canonical
     Tn\emph{2006}. & We did not manually confirm the target-site duplications; relied on
     positional co-location. \\
C5 & Two spatial Tn\emph{7} clusters; TnsA/D/E present; \texttt{dfrA1} and \texttt{aadA1} at both. &
     Did not verify Tn\emph{7} target-site attTn\emph{7}--II duplication or 22\,bp end
     recognition sequences. \\
C6 & Tn\emph{7}$^{+}$ literally the next CDS downstream of \emph{glmS}; \texttt{tet(B)} + \texttt{sul2}
     inside the same $\sim$21.7\,kb block. & Reported span 21.7\,kb differs from paper 22{,}852\,bp
     by $\sim$5\%; we did not reconstruct the paper's exact right-boundary IS\emph{Aba1}. \\
C7 & S83L reproduced exactly; IS\emph{Aba1}--MarR juxtaposition confirmed; second frameshifted
     MarR-family pseudogene bonus. & Transcriptional half untested: no RNAseq, no promoter
     mapping, no functional confirmation that \emph{marA} runs constitutively from an
     IS\emph{Aba1}-internal promoter.  A key half of the mechanism claim rests on inference. \\
\bottomrule
\end{longtable}

\section{GENUINE CRITIQUE of the replication}

This is an honest self-review.  The replication landed on the strong side
of the verdict scale (92), but several aspects should be tempered before
using this result to draw broader conclusions.

\paragraph{We replicated the paper's genome, not the paper's genomics.}
By starting from \texttt{GCA\_021484925.1}, we are analysing the
same assembled sequence the authors submitted.  We therefore validate
that the authors' \emph{annotations} and \emph{interpretations} of that
sequence are reproducible.  We do \textbf{not} validate the upstream
Unicycler assembly itself, and we do not test whether an independent
assembly of the raw MinION+MiSeq reads would recover the same four
small plasmids, the same closure of the chromosome, or the same
resolution of the two Tn\emph{7} copies.  The reads themselves
(SRA / ENA) were not re-fetched or re-assembled.  This is a real gap:
a paper about a \emph{finished long-read assembly} was replicated using
only the finished product.

\paragraph{Position-based claims lean on PGAP.}
Our locus walks (AbaR28, Tn\emph{7} clusters, IS\emph{Aba1}/marR)
rest on NCBI PGAP's feature annotation.  PGAP versions and its
underlying HMMs update; a re-run with a different PGAP release could
re-label some CDSs (particularly pseudogene calls for MarR-family
regulators and \emph{parC}).  We did not lock a specific PGAP release
or re-run \texttt{prokka}/\texttt{bakta} as an independent second
annotator.

\paragraph{Numeric hand-waving on the Tn\emph{7}$^{+}$ span.}
The paper claims 22{,}852\,bp; we report ``$\sim$21.7\,kb''.  We accept
that as ``same order of magnitude, boundary-definition difference'',
but we did not do the disciplined thing --- pinpoint the paper's exact
left/right boundary features, then measure between the same coordinates
in our own annotation.  A more rigorous replication would either
recover 22{,}852\,bp exactly or explicitly identify the missing
$\sim$1.1\,kb.

\paragraph{The C7 transcriptional claim is not truly tested.}
The paper's most \emph{interesting} contribution is a two-step
functional mechanism: (i) gyrA QRDR + (ii) IS\emph{Aba1} inactivates
\emph{marR}, and \emph{marA} is now driven from an IS\emph{Aba1}-internal
promoter.  Step (i) we replicate cleanly.  Step (ii) we replicate
\emph{structurally} (IS\emph{Aba1} is next to a MarR-family regulator;
another MarR-family gene is frameshifted).  We do \emph{not}:
\begin{itemize}[nosep]
  \item Confirm the specific IS\emph{Aba1}--internal promoter is actually
        driving \emph{marA};
  \item Show elevated \emph{marA} transcript in the presence of the
        insertion (no RNAseq);
  \item Show FQ efflux upregulation attributable to \emph{marA}
        constitutivity;
  \item Distinguish among the multiple ISAba1-associated MarR-family loci
        (chr 2{,}317k intact regulator vs.\ chr 1{,}422k frameshifted pseudogene) which is
        actually responsible for the phenotype;
  \item Show that \emph{marR} pseudogene $\to$ FQR phenotype in a
        clean genetic background (knockout / complementation).
\end{itemize}
Scoring the structural half as ``match'' should not be read as
validating the mechanistic model.

\paragraph{Sequence-identity thresholds are permissive.}
AMR calls used \texttt{--minid 90 --mincov 80}.  This is conservative
enough to avoid the worst false positives but permissive enough that
alleles slightly diverged from the reference (e.g.\ variants of
\texttt{aadA1}, promoter mutants of \texttt{bla}$_{\text{OXA-23}}$)
would be lumped in without differentiation.  We did not use e.g.\
\texttt{--minid 98} to stress-test the calls.

\paragraph{LLM judges are correlated evidence, not independent evidence.}
Both judges were given \emph{the same structured dossier we wrote}.
They can only score what we chose to put in front of them.  A 92 from
two LLMs reading our own evidence brief is not two independent
replications --- it is one replication scored twice.

\paragraph{Publication-bias-shaped success.}
This is a paper that describes something the authors were able to
see directly in the finished sequence.  It is almost guaranteed to
replicate against that same sequence; the replication is largely a
check that the annotation is faithful.  A more probing replication
would be against a \emph{different} XDR GC1 isolate --- can we identify
Tn\emph{7}$^{+}$ or the IS\emph{Aba1}--\emph{marR} inactivation pattern
in any other public GC1 genome?  We did not attempt this.

\paragraph{Missing paper PDF.}
The primary source (\texttt{dkac115.pdf}) is not on disk.  All three
open-access endpoints (OUP, Europe PMC, NCBI PMC) failed
non-interactively during the fetch cap; the paper is not in the Eagle
Marker/Nougat central corpus.  Our replication was written against the
abstract, the authors' own GenBank submission metadata, and the paper's
own claims transcribed from the abstract and PMC-cached figures.  Body
text of the paper (Methods, discussion of the transcriptional evidence,
supplementary tables) was not consulted during this pass.  This is the
single largest hole in the audit trail.

\section{Verdict}

\textbf{REPLICATED, score 92/100} (mean of LLM judges 88 and 95).

\emph{Justification.}  Six of seven core claims reproduce
quantitatively from the paper's own public genome using open-source
tools.  C7 reproduces every \emph{structural} component of the
paper's novel FQR-route model; the only untested piece is the
transcriptional-activation step.  No claim is contradicted.  The
paper's headline novel contribution --- the ``Tn\emph{7}$^{+}$'' module
at attTn\emph{7} downstream of \emph{glmS} carrying an AbGRI1-derived
\texttt{tet(B)}/\texttt{sul2} segment --- is confirmed at the predicted
genomic address.

\section{Open Questions}

\paragraph{Q1.}  Does the IS\emph{Aba1}-internal promoter that the
paper posits actually drive \emph{marA} constitutively \emph{in vivo}?
\emph{Basis:} the mechanism is stated as the paper's central novel
contribution but was not tested transcriptionally in the paper (no
RNAseq/qPCR shown for MRSN 56 specifically) and was not tested here.
\emph{Next steps:} obtain a viable MRSN 56 (or reconstructed derivative)
isolate; run RNAseq $\pm$ FQ challenge; compare \emph{marA} transcript
levels to an isogenic \emph{marR}-restored strain; map 5$'$ transcript
ends by RACE / rend-seq to confirm the IS\emph{Aba1}-internal TSS.

\paragraph{Q2.}  Which MarR-family locus is functionally responsible ---
the chr 2{,}317k IS\emph{Aba1}-adjacent \emph{intact} MarR regulator,
the chr 1{,}422k frameshifted MarR pseudogene, or both?
\emph{Basis:} PGAP flags two MarR-family loci with different insults
in MRSN 56; the paper does not explicitly disentangle them.
\emph{Next steps:} construct isogenic knockouts and complementations
of each candidate MarR-family gene in a susceptible GC1 background;
measure FQ MIC and \emph{marA} expression; test which restoration
suppresses the phenotype.

\paragraph{Q3.}  Is the Tn\emph{7}$^{+}$ configuration
(attTn\emph{7}--Tn\emph{7}--AbGRI1-tetB/sul2 fusion) unique to MRSN 56,
or has it disseminated across GC1 / XDR \emph{A.~baumannii} more
widely?  \emph{Basis:} the paper reports it as a novel structural
event but does not survey other genomes.  \emph{Next steps:} run a
structural-variant scan (nucmer / minimap2) across all $\sim$16{,}000
\emph{A.~baumannii} genomes in NCBI RefSeq for the diagnostic
Tn\emph{7}$^{+}$ junction (\emph{glmS}--TnsA--\dots--\emph{tet(B)}--
\emph{sul2}--IS\emph{Aba1}); phylogenetically cluster hits and infer
whether Tn\emph{7}$^{+}$ has moved horizontally.

\paragraph{Q4.}  How reproducible is this MinION+MiSeq Unicycler
assembly at the finished level?  \emph{Basis:} the paper's finished
genome is the substrate for every downstream claim; we replicated
against the finished product, not the reads.  \emph{Next steps:} pull
raw reads from ENA/SRA (BioProject PRJNA742487, biosample SAMN20178847);
re-run Unicycler v0.4.0 and modern hybrid assemblers (Trycycler,
Flye+polypolish); check whether the four small plasmids and the two
Tn\emph{7} copies are recovered independently, and whether the
Tn\emph{7}$^{+}$ configuration is boundary-stable under different
assemblers.

\paragraph{Q5.}  What are the fitness costs / clinical stability of the
IS\emph{Aba1}$\to$\emph{marR} inactivation route to FQ resistance
compared with the classical \emph{gyrA}+\emph{parC} double-mutant
route?  \emph{Basis:} the paper proposes a novel two-step \emph{route}
but does not compare its stability, its reversibility, or its
downstream MIC ceiling with the more familiar gyrA+parC pathway (which
MRSN 56 also carries, in the form of the \emph{parC} pseudogene we
flagged as a bonus finding).  \emph{Next steps:} in a common GC1
background, construct four strains
--- (a) WT, (b) gyrA S83L only, (c) IS\emph{Aba1}$\to$\emph{marR} only,
(d) both --- and measure FQ MIC, growth rate, competitive fitness,
and reversion rate over serial passage; ask whether the
IS\emph{Aba1}-route strains bear a lower fitness cost and are therefore
more likely to persist clinically.

\section{Provenance}

Assembly: \texttt{GCA\_021484925.1} (chromosome CP090606.1, plasmids
CP080453--CP080456).  Reference protein: \texttt{WP\_000116449.1}.
Environment: \texttt{uicgpu:/data/stevens/envs/bvbrc14}.
LLM proxy: \texttt{argo:localhost:44497} (models
\texttt{argo:gpt-5.1} and \texttt{argo:gemini-2.5-pro}, free ANL endpoint).
No paid API used.

\end{document}
